% !TEX root = ../main.tex

\chapter{آزمایش‌ها و نتایج}

% =========================================================
% 6-1
% =========================================================
\section{مقدمه}

در این فصل، روش انجام آزمایش‌ها و نتایج حاصل از ارزیابی
مؤلفه‌های اصلی سامانه ارائه می‌شوند.

آزمایش‌ها در سه بخش اصلی انجام شده‌اند:

\begin{itemize}
    \item ارزیابی مدل‌های تشخیص گفتار؛
    \item ارزیابی مدل‌های زبانی؛
    \item ارزیابی مدل‌های تبدیل متن به گفتار.
\end{itemize}

هدف، انتخاب مدلی نیست که تنها در یک معیار بهترین عملکرد را
داشته باشد، بلکه تلاش شده است توازن میان دقت، کیفیت،
زمان اجرا و مصرف منابع بررسی شود.

تحلیل دلایل تفاوت میان مدل‌ها و انتخاب نهایی در فصل بعد
ارائه خواهد شد.


% =========================================================
% 6-2
% =========================================================
\section{ارزیابی مدل‌های تشخیص گفتار}

برای ارزیابی مدل‌های تشخیص گفتار، از بخش آزمون
\LTRfootnote{Test Split}
مجموعه‌داده
\lr{Persian Common Voice 25}
استفاده شد.

این بخش شامل
\lr{10,519}
نمونه صوتی با مجموع مدت زمان

\begin{center}
\lr{52,047.28 seconds}
\end{center}

معادل تقریباً 14.46 ساعت گفتار است.

چهار مدل زیر تحت شرایط ارزیابی یکسان مورد آزمایش قرار گرفتند:

\begin{itemize}

    \item \lr{wav2vec2-large-xlsr-persian-v3}

    \item \lr{Whisper-large-v3}

    \item \lr{vosk-model-fa-0.42}

    \item \lr{stt\_fa\_fastconformer\_hybrid\_large}

\end{itemize}

برای مقایسه مدل‌ها از معیارهای
\lr{WER}،
\lr{CER}،
\lr{RTF}،
متوسط زمان تأخیر،
حافظه مصرفی \lr{GPU}
و مجموع زمان استنتاج استفاده شد.


% ---------------------------------------------------------
\subsection{نتایج ارزیابی}

\begin{table}[H]
\centering
\caption{نتایج ارزیابی مدل‌های تشخیص گفتار روی بخش آزمون
\lr{Persian Common Voice 25}}
\label{tab:asr_results}

\resizebox{\textwidth}{!}{
\begin{tabular}{lrrrrrrr}
\toprule

\textbf{\lr{Model}} &
\textbf{Samples} &
\textbf{\lr{WER}} &
\textbf{\lr{CER}} &
\textbf{\lr{RTF}} &
\textbf{\lr{Latency} (s)} &
\textbf{\lr{GPU} (GB)} &
\textbf{Inference (s)}
\\

\midrule

\lr{wav2vec2-large-xlsr-persian-v3}
&
10,519
&
0.2615
&
0.1205
&
\textbf{0.0150}
&
\textbf{0.0744}
&
1.64
&
\textbf{782.60}
\\

\lr{Whisper-large-v3}
&
10,519
&
0.3677
&
0.1064
&
0.2944
&
1.4568
&
6.70
&
15,324.19
\\

\lr{vosk-model-fa-0.42}
&
10,519
&
0.2234
&
0.0593
&
0.3051
&
1.5098
&
\textbf{0.00}
&
15,881.76
\\

\lr{stt\_fa\_fastconformer\_hybrid\_large}
&
10,519
&
\textbf{0.1233}
&
\textbf{0.0379}
&
0.0275
&
0.1359
&
0.88
&
1,429.43
\\

\bottomrule
\end{tabular}
}
\end{table}

نتایج نشان می‌دهند که مدل
\lr{stt\_fa\_fastconformer\_hybrid\_large}
بهترین دقت را در میان مدل‌های مورد بررسی ارائه کرده است.
این مدل به
\lr{WER}
برابر با
$0.1233$
و
\lr{CER}
برابر با
$0.0379$
دست یافت که کمترین مقدار هر دو معیار در این آزمایش است.

از نظر سرعت،
\lr{wav2vec2-large-xlsr-persian-v3}
بهترین عملکرد را ارائه کرد.
این مدل دارای
\lr{RTF}
برابر با
$0.0150$
و متوسط تأخیر
$0.0744$
ثانیه برای هر نمونه بود.
مجموع زمان استنتاج آن نیز تنها
$782.60$
ثانیه برای کل مجموعه آزمون بود.

مدل
\lr{FastConformer Hybrid}
با وجود آنکه از \lr{wav2vec2} کندتر بود، همچنان سرعت بسیار مناسبی
داشت.
\lr{RTF} این مدل برابر با
$0.0275$
و متوسط زمان پردازش هر نمونه
$0.1359$
ثانیه بود.
بنابراین این مدل علاوه بر بهترین دقت، از نظر سرعت نیز برای
یک کاربرد تعاملی گزینه مناسبی محسوب می‌شود.

مدل
\lr{Vosk}
بدون استفاده از حافظه \lr{GPU} اجرا شد که مهم‌ترین مزیت آن از
دیدگاه محدودیت منابع است.
با این حال، \lr{RTF} برابر با
$0.3051$
و متوسط تأخیر
$1.5098$
ثانیه نشان می‌دهد که اجرای آن در تنظیمات این آزمایش نسبت به
\lr{wav2vec2} و \lr{FastConformer} کندتر بوده است.

مدل
\lr{Whisper Large-v3}
بیشترین \lr{WER} را با مقدار
$0.3677$
و بیشترین مصرف حافظه \lr{GPU} را با مقدار تقریبی
$6.70$
گیگابایت نشان داد.
همچنین \lr{RTF} آن برابر با
$0.2944$
بود که به‌طور قابل توجهی بیشتر از \lr{wav2vec2} و \lr{FastConformer}
است.

در مجموع، نتایج این آزمایش نشان می‌دهند که
\lr{FastConformer Hybrid} بهترین توازن میان دقت، سرعت و مصرف
منابع را در مجموعه آزمون مورد استفاده ارائه کرده است.
% =========================================================
% 6-3
% =========================================================
\section{ارزیابی مدل‌های زبانی}

ارزیابی اصلی مدل‌های زبانی با استفاده از دو وظیفه از مجموعه‌داده
\lr{ParsiNLU} انجام شد: پرسش‌های چندگزینه‌ای و درک مطلب
\lr{ParsBench RC}. همه مدل‌ها با تنظیمات یکسان کوانتیزه‌سازی
چهاربیتی \lr{NF4} اجرا شدند تا مقایسه تا حد امکان منصفانه باشد.

مدل‌های ارزیابی‌شده عبارت‌اند از:

\begin{itemize}
    \item \lr{gemma-2-2b-it}
    \item \lr{gemma-2-9b-it}
    \item \lr{Qwen2.5-3B-Instruct}
    \item \lr{Qwen2.5-7B-Instruct}
    \item \lr{Dorna2-Llama3.1-8B-Instruct}
\end{itemize}

این مدل‌ها پس از یک غربال اولیه استقراری انتخاب شدند: هر مدل باید
در حالت چهاربیتی \lr{NF4} به‌طور کامل در حافظه \lr{GPU} محلی پروژه
قرار می‌گرفت. در نتیجه، مدل‌های بزرگ‌تری که برای بارگذاری به چند
شتاب‌دهنده، انتقال وزن‌ها به حافظه اصلی یا سرویس ابری نیاز داشتند،
خارج از دامنه مقایسه بودند.
در نتیجه، رتبه‌بندی گزارش‌شده عملکرد نسخه‌های چهاربیتی را نشان
می‌دهد و نباید بدون آزمایش جداگانه به نسخه‌های تمام‌دقت تعمیم داده شود.

وظیفه چندگزینه‌ای شامل \lr{1050} پرسش بود که به‌طور مساوی در سه
دسته دانش عمومی، ادبیات و ریاضی و منطق توزیع شدند؛ بنابراین هر
دسته شامل \lr{350} پرسش بود. معیار اصلی در این وظیفه دقت است.

بخش درک مطلب \lr{ParsBench RC} شامل \lr{600} نمونه با فیلدهای
\lr{instruction}، \lr{input} و \lr{output} بود. خروجی مدل‌ها در این
وظیفه با معیارهای تطابق دقیق، \lr{F1} در سطح توکن و شباهت معنایی
ارزیابی شد.


% ---------------------------------------------------------
\subsection{نتایج پرسش‌های چندگزینه‌ای}

\begin{table}[H]
\centering
\caption{دقت مدل‌های زبانی در وظیفه چندگزینه‌ای \lr{ParsiNLU}}
\label{tab:llm_mcq_results}

\resizebox{\textwidth}{!}{
\begin{tabular}{lcccc}
\toprule
\textbf{\lr{Model}} &
\textbf{دانش عمومی} &
\textbf{ادبیات} &
\textbf{ریاضی و منطق} &
\textbf{کل} \\
\midrule

\lr{Gemma 2 9B} &
\textbf{56.57\%} &
\textbf{47.71\%} &
36.29\% &
\textbf{46.86\%} \\

\lr{Qwen2.5 7B} &
45.14\% &
33.43\% &
\textbf{42.57\%} &
40.38\% \\

\lr{Dorna2 8B} &
40.86\% &
37.43\% &
32.57\% &
36.95\% \\

\lr{Qwen2.5 3B} &
33.71\% &
26.86\% &
28.00\% &
29.52\% \\

\lr{Gemma 2 2B} &
34.29\% &
25.43\% &
17.43\% &
25.71\% \\

\bottomrule
\end{tabular}
}
\end{table}

\lr{Gemma 2 9B} با دقت کل \lr{46.86\%} رتبه نخست را به دست آورد
و در دو دسته دانش عمومی و ادبیات نیز بهترین مدل بود. مدل
\lr{Qwen2.5 7B} با دقت کل \lr{40.38\%} در رتبه دوم قرار گرفت، اما
با دقت \lr{42.57\%} بهترین عملکرد را در ریاضی و منطق نشان داد.
پس از آن، \lr{Dorna2 8B}، \lr{Qwen2.5 3B} و \lr{Gemma 2 2B} به‌ترتیب
در رتبه‌های بعدی قرار گرفتند.

\subsection{ارزیابی درک مطلب}

در وظیفه درک مطلب، همان پنج مدل و همان تنظیمات چهاربیتی \lr{NF4}
روی هر \lr{600} نمونه \lr{ParsBench RC} ارزیابی شدند. استفاده هم‌زمان
از تطابق دقیق، \lr{F1} توکنی و شباهت معنایی امکان بررسی پاسخ‌ها را
از سه منظر مکمل فراهم کرد: انطباق کامل با پاسخ مرجع، همپوشانی
واژگانی و نزدیکی معنایی. این طراحی نسبت به اتکا به یک معیار منفرد،
برای پاسخ‌های تولیدی فارسی ارزیابی دقیق‌تری فراهم می‌کند.

\subsection{آزمایش مقدماتی \lr{FarsInstruct}}

پیش از این ارزیابی‌ها، آزمایشی مقدماتی روی \lr{1000} مثال تصادفی
از \lr{FarsInstruct} انجام شد. به‌دلیل
ترکیب تصادفی دسته‌ها و کنترل‌نشدن سهم هر نوع وظیفه، این آزمایش
صرفاً به‌عنوان مطالعه راهنما در نظر گرفته می‌شود و مبنای رتبه‌بندی
اصلی مدل‌ها نیست. نتایج آن برای طراحی ارزیابی‌های بعدی و انتخاب
معیارها استفاده شد؛ در حالی که نتیجه‌گیری اصلی این گزارش بر وظایف
مشخص و کنترل‌شده \lr{ParsiNLU} و \lr{ParsBench RC} استوار است.


% =========================================================
% 6-4
% =========================================================
\section{ارزیابی مدل‌های تبدیل متن به گفتار}

سه مدل \lr{TTS} فارسی مورد آزمایش قرار گرفتند:

\begin{itemize}
    \item \lr{Kamtera\_VITS}
    \item \lr{VITS\_MMS\_FAS}
    \item \lr{PIPER\_TTS\_FAS}
\end{itemize}

ارزیابی روی
\lr{1052}
نمونه انجام شد.

برای بررسی کیفیت ادراکی از \lr{NISQA} و برای بررسی
قابل‌فهم‌بودن گفتار از \lr{WER} و \lr{CER} مبتنی بر بازشناسی مجدد
صوت استفاده شد.


% ---------------------------------------------------------
\subsection{نتایج کیفیت ادراکی}

\begin{table}[H]
\centering
\caption{نتایج \lr{NISQA} برای مدل‌های \lr{TTS}}
\label{tab:tts_nisqa}

\resizebox{\textwidth}{!}{
\begin{tabular}{lccccc}
\toprule
\textbf{\lr{Model}} &
\textbf{\lr{MOS}} &
\textbf{\lr{Noise}} &
\textbf{\lr{Discontinuity}} &
\textbf{\lr{Coloration}} &
\textbf{\lr{Loudness}} \\
\midrule

\lr{Kamtera VITS} &
3.380 &
3.632 &
4.077 &
3.814 &
4.098 \\

\lr{MMS VITS} &
3.913 &
4.011 &
4.048 &
3.633 &
3.887 \\

\lr{Piper} &
4.534 &
4.070 &
4.656 &
4.168 &
4.282 \\

\bottomrule
\end{tabular}
}
\end{table}

\lr{Piper} بالاترین امتیاز کلی \lr{NISQA} را با مقدار
$4.534$
به دست آورده است.
این مدل همچنین در معیارهای ناپیوستگی، coloration و loudness
عملکرد بالایی داشته است.


% ---------------------------------------------------------
\subsection{نتایج قابل فهم بودن و سرعت}

\begin{table}[H]
\centering
\caption{نتایج دقت و سرعت مدل‌های \lr{TTS}}
\label{tab:tts_runtime}

\resizebox{\textwidth}{!}{
\begin{tabular}{lccccc}
\toprule
\textbf{\lr{Model}} &
\textbf{\lr{WER}} &
\textbf{\lr{CER}} &
\textbf{\lr{RTF}} &
\textbf{\lr{Latency} (s)} &
\textbf{\lr{Speedup}} \\
\midrule

\lr{Kamtera VITS} &
0.144 &
0.068 &
0.043 &
0.175 &
23.27x \\

\lr{MMS VITS} &
0.253 &
0.114 &
0.032 &
0.106 &
31.55x \\

\lr{Piper} &
0.107 &
0.045 &
0.192 &
0.948 &
5.20x \\

\bottomrule
\end{tabular}
}
\end{table}

مدل \lr{MMS} سریع‌ترین مدل \lr{TTS} در این آزمایش بوده است و
کمترین \lr{RTF} را به دست آورده است.

در مقابل، \lr{Piper} با وجود سرعت کمتر، کمترین \lr{WER} و \lr{CER} و
بالاترین امتیاز کیفیت \lr{NISQA} را ارائه کرده است.

بنابراین نتایج نشان‌دهنده یک توازن روشن میان کیفیت و سرعت
در انتخاب مدل \lr{TTS} هستند.
% =========================================================
% 6-7
% =========================================================
\section{جمع‌بندی}

در این فصل، مدل‌های مختلف \lr{ASR}، \lr{LLM} و \lr{TTS} به‌صورت مستقل
ارزیابی شدند.

در بخش تشخیص گفتار، \lr{FastConformer Hybrid} بهترین دقت را
ارائه کرد و در عین حال سرعت مناسبی داشت.

در بخش مدل زبانی، \lr{Gemma 2 9B} با دقت کل \lr{46.86\%} بهترین
عملکرد را در وظیفه چندگزینه‌ای \lr{ParsiNLU} ارائه کرد؛ در حالی که
\lr{Qwen2.5 7B} در دسته ریاضی و منطق بهترین نتیجه را به دست آورد.

در بخش تبدیل متن به گفتار نیز \lr{Piper} بالاترین کیفیت ادراکی و
بهترین قابل‌فهم‌بودن خروجی را به دست آورد، هرچند \lr{MMS} از نظر
سرعت تولید بهتر بود.

با وجود برتری دقت \lr{Gemma 2 9B}، این نتیجه به‌تنهایی برای جایگزینی
مدل سامانه نهایی کافی نیست؛ زیرا ارقام به‌روزشده زمان اجرا و مصرف
حافظه برای هر پنج مدل در این مقایسه گزارش نشده‌اند. ازاین‌رو
\lr{Gemma 2 2B} همچنان کاندید اجرای محلی کم‌منبع باقی می‌ماند و
\lr{Gemma 2 9B} کاندید مبتنی بر بیشترین دقت است.

با این حال، انتخاب نهایی تنها پس از بررسی عملکرد کل سامانه،
تأخیر انتهابه‌انتها و مصرف منابع در شرایط اجرای واقعی انجام
خواهد شد.

در فصل بعد، دلایل تفاوت عملکرد مدل‌ها، توازن میان دقت و سرعت،
محدودیت‌های آزمایش‌ها و پیامدهای نتایج برای طراحی سامانه
مورد بحث قرار خواهند گرفت.
